\documentclass[12pt,a4paper]{article}

% Pacotes necessários
\usepackage[utf8]{inputenc}
\usepackage[T1]{fontenc}
\usepackage[portuguese]{babel}
\usepackage{geometry}
\usepackage{listings}
\usepackage{xcolor}
\usepackage{graphicx}
\usepackage{hyperref}
\usepackage{amsmath}
\usepackage{booktabs}
\usepackage{enumitem}

% Configuração de página
\geometry{margin=2.5cm}

% Configuração de cores para código
\definecolor{codegreen}{rgb}{0,0.6,0}
\definecolor{codegray}{rgb}{0.5,0.5,0.5}
\definecolor{codepurple}{rgb}{0.58,0,0.82}
\definecolor{backcolour}{rgb}{0.95,0.95,0.92}

% Configuração do estilo de código C#
\lstdefinestyle{csharp}{
    backgroundcolor=\color{backcolour},
    commentstyle=\color{codegreen},
    keywordstyle=\color{blue},
    numberstyle=\tiny\color{codegray},
    stringstyle=\color{codepurple},
    basicstyle=\ttfamily\footnotesize,
    breakatwhitespace=false,
    breaklines=true,
    captionpos=b,
    keepspaces=true,
    numbers=left,
    numbersep=5pt,
    showspaces=false,
    showstringspaces=false,
    showtabs=false,
    tabsize=2,
    language=[Sharp]C,
    frame=single
}

% Configuração do estilo para YAML
\lstdefinestyle{yaml}{
    backgroundcolor=\color{backcolour},
    commentstyle=\color{codegreen},
    keywordstyle=\color{blue},
    basicstyle=\ttfamily\footnotesize,
    breaklines=true,
    frame=single
}

\lstset{style=csharp}

% Configuração de hyperlinks
\hypersetup{
    colorlinks=true,
    linkcolor=blue,
    filecolor=magenta,
    urlcolor=cyan,
}

\begin{document}

\title{\textbf{Sistema RPC Distribuído com RabbitMQ}\\
\large{Implementação do Cliente e Servidor}}
\author{Vinícius}
\date{\today}

\maketitle

\newpage

\section{Implementação do Cliente e Servidor}

O sistema foi desenvolvido em C\# com duas aplicações .NET independentes: servidor (.NET 8.0) e cliente (.NET 9.0), comunicando-se exclusivamente através de filas RabbitMQ.

\subsection{Servidor}

O servidor utiliza o padrão \textit{Strategy} através da interface \texttt{IOperationService}, permitindo extensibilidade através do registro de serviços em um \texttt{Dictionary}. Implementa reconexão automática (até 10 tentativas) e opera em duas filas: \texttt{fila\_rpc} (operações síncronas) e \texttt{fila\_async} (fire-and-forget). O processamento utiliza \texttt{EventingBasicConsumer} com confirmação manual via \texttt{BasicAck}/\texttt{BasicNack}.

\subsection{Cliente}

O cliente implementa \texttt{IDisposable} para gerenciamento de conexões e expõe dois métodos: \texttt{Call()} para RPC síncrono (com \texttt{CorrelationId}, fila temporária e timeout de 5s) e \texttt{SendAsync()} para envio assíncrono sem resposta. As mensagens são serializadas em JSON através da classe \texttt{RequestMessage}.

\section{Bibliotecas Utilizadas}

\begin{itemize}
    \item \textbf{RabbitMQ.Client 6.8.1}: Biblioteca oficial para comunicação AMQP (ConnectionFactory, IModel, EventingBasicConsumer)
    \item \textbf{System.Text.Json}: Serialização/desserialização nativa do .NET
    \item \textbf{System.IO}: Operações de arquivo (FileService)
    \item \textbf{System.Threading}: Gerenciamento de timeout e delays
\end{itemize}

\section{Configuração do Ambiente}

O ambiente utiliza Docker Compose com três serviços: RabbitMQ (portas 5672 e 15672, com healthcheck), servidor containerizado com volume para persistência (\texttt{./server:/data}), e rede bridge dedicada. Variáveis de ambiente configuram hostname do broker, nomes das filas e timeout RPC.

\section{Operações Implementadas}

Todas as operações implementam a interface \texttt{IOperationService} com método \texttt{ExecuteAsync()}.

\subsection{MessageService (msg)}
Recebe mensagem de texto e retorna confirmação com timestamp e contagem de caracteres. Implementa validação de entrada vazia.

\subsection{FileService (file)}
Escreve conteúdo em \texttt{/data/file.txt} usando \texttt{AppendAllTextAsync()} para I/O assíncrono. Adiciona timestamp automático e retorna tamanho total do arquivo. Arquivo persiste via volume Docker.

\subsection{MathService (calc)}
Executa 7 operações matemáticas (soma, subtração, multiplicação, divisão, potência, módulo, raiz n-ésima) com formato \texttt{operacao,num1,num2}. Implementa validação de formato, parsing com \texttt{double.TryParse()}, tratamento de divisão/módulo por zero e raiz com índice zero. Aceita múltiplos aliases por operação (ex: soma/+).

\end{document}
